\documentclass[aps,pre,reprint,nofootinbib]{revtex4-2}

\usepackage{amsmath,amssymb,bm}
\usepackage{array}
\usepackage{booktabs}
\newcolumntype{L}[1]{>{\raggedright\arraybackslash}p{#1}}
\usepackage{graphicx}
\usepackage{hyperref}
\usepackage{microtype}
\usepackage{pgfplots}
\usepgfplotslibrary{groupplots}
\pgfplotsset{compat=1.15,
  paperaxis/.style={
    width=\columnwidth, height=0.58\columnwidth,
    grid=major, grid style={gray!20},
    tick label style={font=\footnotesize},
    label style={font=\footnotesize},
    legend style={font=\footnotesize, draw=none, fill=none},
    line width=0.9pt}}

\newcommand{\dd}{\mathrm{d}}

\begin{document}

\title{Three Sources of Observed Entropy in Reversible Hidden Dynamics:\\
Mixing, Projection Multiplicity, and Inaccessible Degrees of Freedom}

\author{Andrew H. Bond}
\email{andrew.bond@sjsu.edu}
\thanks{ORCID: 0009-0003-2599-6158}
\affiliation{San Jos\'e State University, San Jos\'e, California, USA}
\date{August 5, 2026}

\begin{abstract}
A previous paper showed that folds of a singular projection create observed
entropy with zero hidden information loss.  This paper completes that
program with three measured experiments that say where observed entropy
comes from and when it becomes real.  The first experiment is a matched
design in which an integrable and a chaotic dynamics share every other
parameter while a folding observer and a monotone observer read the same
trajectories.  Sustained growth of the observed ensemble entropy follows
the dynamics under either observer, the multiplicity of hidden states
behind one observation follows the observer under either dynamics, and the
two mechanisms are independent.  The second experiment measures a hierarchy
of consumers reading a double fold with increasing access.  Their exact
conditional entropies form a lawful decreasing chain whose gaps price each
additional observable, and the chain is the classical counterpart of a
quantum consumer hierarchy measured earlier in the campaign.  The third
experiment couples the earlier reversible fold cycle to a thermal bath the
observer cannot reverse.  Reversal restricted to the system fails in
proportion to the coupling, reversal extended to the bath succeeds exactly
at every coupling, and the record of the leak is invisible to the coarse
bath observables measured here.  Within this model class, observed entropy
growth is mixing, observed ambiguity is the observer's own singularities,
and irreversibility is an accounting of which degrees of freedom the
observer can reach.  Nothing in this paper is a claim about physical
spacetime or about the thermodynamics of any real system.
\end{abstract}

\maketitle

\section{Introduction}
\label{sec:intro}

An observer watching a hidden dynamical system through a projection can see
its recorded entropy change for reasons that have nothing to do with each
other.  The hidden flow can genuinely mix, spreading an ensemble across the
observer's bins.  The projection can fold, so that a single observed
instant sits over several hidden states at once.  And information can leak
into degrees of freedom the observer cannot manipulate, so that a reversal
protocol available to the observer no longer undoes the evolution.  These
three sources are routinely described with the same word, entropy, and a
central discipline of the present campaign is to refuse that conflation
until an experiment forces it or forbids it.

The first paper of this track \cite{EntropyPaper} studied the fold itself.
It proved that a nondegenerate fold changes the number of hidden preimages
by two while conserving an orientation-weighted count, computed exact
branch entropies for canonical control maps, measured the resolution
scaling of the caustic, and demonstrated a reversible Hamiltonian ensemble
whose observed entropy swings by more than two bits and then returns every
bit, with the initial hidden ensemble recovered to near machine precision.
That paper established that folds create observed entropy without hidden
loss.  It deliberately left three questions open.  Whether fold
multiplicity or hidden mixing controls the growth of observed entropy in a
system that has both.  What a hierarchy of observers with graded access is
worth in exact conditional bits.  And when the recoverable ambiguity of a
projection becomes irreversible entropy production.

This paper answers all three within the same model family, and the track
is thereby complete.  The answers have a common shape.  Each candidate
source of observed entropy is isolated by a matched design in which the
other sources are held fixed, and each answer is a measured number with a
falsification bar that was written down before the run.  The mixing
experiment finds that growth belongs to the dynamics and multiplicity
belongs to the observer, with no leakage between the two.  The hierarchy
experiment finds a lawful chain of exact conditional entropies whose gaps
price orientation and branch knowledge separately.  The bath experiment
finds that irreversibility is not a property of the dynamics at all in
this class.  It is a property of which momenta the observer is permitted
to flip, and the leaked record hides in correlations that no coarse bath
observable measured here can see.

Every result is a
computation or a measured simulation in a declared toy model class, graded
as measured in model at the exploratory label, the grade for measured
simulations whose evidence records carry the exploratory label because no
seal governs them.  No claim is
made about physical spacetime, about the thermodynamics of any real
system, or about quantum gravity.

\section{The model family and the completed track}
\label{sec:family}

All three experiments use the campaign's standard hidden system, two
coupled degrees of freedom with Hamiltonian
\begin{equation}
H = \tfrac{1}{2}\left(p_t^2 + p_u^2\right)
  + \tfrac{1}{2}\omega_t^2 t^2 + \tfrac{1}{2}\omega_u^2 u^2
  + \tfrac{\lambda}{4} u^4 + gE\,t u,
\label{eq:toy}
\end{equation}
with $\omega_u = 1.2$ and $g = 0.25$ throughout, integrated by
time-symmetric methods whose energy drift is monitored in every run.  The
observed coordinate is $t$ and the observer's records are binned at a
declared resolution $\epsilon$.  The projection of interest sends the
hidden evolution parameter $\tau$ to an observed time functional, and the
observer's singularities are the critical points of that functional.

The track now comprises six experiments and one theorem.  The fold
theorem and the exact
branch entropies, the caustic resolution scaling, and the reversible fold
cycle are published \cite{EntropyPaper}.  The three completions reported
here are the mixing comparison, the consumer hierarchy, and the bath
coupling.  Every result-bearing run writes an append-only evidence record
with hashes, and Sec.~\ref{sec:methods} lists them.

\section{Mixing versus multiplicity}
\label{sec:pe3}

\subsection{A matched two-by-two}

The design question is whether the growth of observed entropy is
controlled by hidden chaotic mixing or by the multiplicity of the
observer's projection.  A four-system comparison invites leakage, because
a system built to mix usually also folds differently.  The experiment
instead crosses two dynamics with two observers so that neither axis can
contaminate the other.

The two dynamics share every parameter of Eq.~(\ref{eq:toy}) with
$\omega_t = 1$ and $E = 2$ and differ only in the quartic coefficient,
$\lambda = 0$ for the integrable member and $\lambda = 0.35$ for the
chaotic member.  The two observers read the same trajectories through two
time functionals, the folding functional $T = t$ and the monotone
functional $T = t + 3\tau$.  The monotone functional never folds because
$\dd T/\dd\tau = p_t + 3$ and $|p_t|$ stays well below $3$ in this family.
The decisive property of the design is that the two functionals differ
only by a linear term in $\tau$, so at every instant the ensemble marginal
of $t$ is the same for both observers.  Growth of the binned ensemble
entropy therefore cannot depend on the observer, and the multiplicity of
worldline states behind one observed level cannot depend on anything but
the observer.  The two mechanisms are separated by construction, and the
experiment measures whether each one actually moves.

A $20000$-member ensemble with thermal hidden initial conditions is
integrated to $\tau = 40$ for each dynamics, and the centered binned
entropy of the observed marginal is recorded along the way at
$\epsilon = 0.05$.  On a reference trajectory the unsigned crossing count
of each observed level, the multiplicity a worldline observer must
resolve, is measured for both functionals.

\subsection{Result}

\begin{figure}[t]
\begin{tikzpicture}
\begin{axis}[paperaxis,
  xlabel={hidden evolution parameter $\tau$},
  ylabel={observed entropy [bits]},
  xmin=0, xmax=40, ymin=0,
  legend pos=south east]
\addplot[blue!70!black, thick]
  table[x=tau, y=integrable] {figdata/pe3_curves.dat};
\addlegendentry{integrable, $\lambda=0$}
\addplot[red!70!black, thick, densely dashed]
  table[x=tau, y=chaotic] {figdata/pe3_curves.dat};
\addlegendentry{chaotic, $\lambda=0.35$}
\end{axis}
\end{tikzpicture}
\caption{Observed binned entropy of the $t$ marginal for the two matched
dynamics.  The early rise is common linear dispersion and does not
separate the systems.  The late behavior does.  Over the second half of
the run the integrable curve recurs with fitted slope $-0.024$ bits per
unit $\tau$ while the chaotic curve continues to grow at $+0.030$ bits
per unit $\tau$.  Data from the committed evidence record of the mixing
experiment.}
\label{fig:pe3}
\end{figure}

Figure~\ref{fig:pe3} shows the two entropy curves.  The total rise over
the run is nearly identical, $5.71$ bits for the integrable dynamics and
$5.83$ bits for the chaotic one, because the early spread of the ensemble
is ordinary linear dispersion and is common to both.  The total rise
therefore does not discriminate.  The late-time
slope discriminates.  Fitted over the second half of the run, the
integrable curve turns over and recurs at $-0.024$ bits per unit $\tau$
while the chaotic curve keeps growing at $+0.030$ bits per unit $\tau$.
Sustained growth follows the dynamics, under either observer, because the
observers share the marginal by construction.

Multiplicity behaves in the complementary way.  The folding observer sees
a mean unsigned crossing multiplicity between $9.6$ and $10.3$ per
observed level, with a maximum of $12$, under both dynamics.  The
monotone observer sees exactly one crossing per level everywhere, under
both dynamics.  Multiplicity follows the observer's singularities and is
indifferent to whether the hidden flow mixes.

The falsification bar, fixed before the run, was met.  A chaotic late
slope at or below the integrable one, a folding observer with mean
multiplicity at or below $1.5$, or a monotone observer seeing any level
more than once would each have invalidated the reading.  The conclusion
is that mixing and multiplicity are independent axes.  Mixing pumps
observed entropy.  Folds create observed ambiguity.  Neither implies the
other.

\section{The consumer hierarchy}
\label{sec:pe5}

\subsection{Graded access to a double fold}

The second completion prices what each additional observable is worth to
an observer of a singular projection.  The system is the double fold
$t = \tau^3 - \tau$ on the declared interval $|\tau| \le 1.6$, the map
whose three branches carry only two orientations because the outer
branches share the sign of $\dd t/\dd\tau$.  That degeneracy is the reason
this map is the right one.  A consumer who learns the orientation has not
learned the branch, so the rungs of the hierarchy are genuinely distinct.

The hidden state is uniform on a $\tau$ grid of spacing $10^{-4}$ and the
observed time is binned at $\epsilon = 0.05$.  Three consumers read the
system.  The first sees only the binned observed time.  The second sees
the binned time and the orientation, the sign of $\dd t/\dd\tau$.  The
third sees the binned time and the full branch label.  Because the state
space is a finite grid, every conditional entropy is an exact finite sum
with no estimator and no sampling error \cite{CoverThomas2006}.

\subsection{Result}

\begin{table}[t]
\caption{The measured consumer chain on the double fold.  Each row states
what the consumer reads and the exact conditional entropy of the hidden
state given that reading.  The residual of the last rung is within-branch
positional uncertainty set by the declared grid and bin resolutions.}
\label{tab:pe5}
\begin{ruledtabular}
\begin{tabular}{lc}
Consumer reads & $H(\text{hidden}\,|\,\text{reading})$ [bits] \\
\colrule
binned position only          & $9.5555$ \\
position and orientation      & $8.8307$ \\
position and branch label     & $8.5453$ \\
full hidden state             & $0$ \\
\end{tabular}
\end{ruledtabular}
\end{table}

Table~\ref{tab:pe5} shows the chain.  The ordering
$9.5555 \ge 8.8307 \ge 8.5453 \ge 0$ holds with strictly positive gaps at
every rung.  Orientation is worth $0.7248$ bits.  The branch label is
worth a further $0.2854$ bits beyond orientation, and that gap is exactly
the price of the outer-branch degeneracy, since orientation resolves the
middle branch but cannot split the outer pair.  The residual $8.5453$ bits
of the branch consumer is within-branch positional uncertainty fixed by
the grid and bin resolutions, not by the fold.  A branch is an interval of
hidden states rather than a point, so no discrete label can take the chain
to zero.

The campaign's quantum track measured the same structure in an Ising
consumer hierarchy, where coarser consumers pay quantifiable
distinguishability costs in a lawful degradation chain
\cite{QuantumPaper}.  The present chain is the classical counterpart that
the quantum paper's reduction argument cites.  The gaps are classical
prices in exact bits, and the ordering is forced by the data-processing
inequality in its most elementary form.

\section{Reversal reach and entropy production}
\label{sec:pe4}

\subsection{The cycle with a bath attached}

The published reversible fold cycle showed that more than two bits of
observed entropy can appear and vanish with zero hidden loss, because
flipping every momentum retraces the evolution exactly
\cite{EntropyPaper}.  The third completion asks what changes when the
system is coupled to an environment.  The hidden coordinate $u$ of
Eq.~(\ref{eq:toy}), with $\omega_t = 1$, $E = 1$, and $\lambda = 0.1$, is
coupled to a bath of $32$ harmonic modes with frequencies spread uniformly
over $[0.5, 2.5]$, thermal initial conditions, and bilinear coupling of
strength $\kappa/\sqrt{32}$ per mode in the style of the standard
oscillator-bath models \cite{CaldeiraLeggett1983}.  The full evolution
remains Hamiltonian and time-symmetric, so in the full space nothing is
irreversible at any coupling.

The observer's limitation is operational.  The observer can flip the
system momenta but cannot touch the bath.  For each coupling
$\kappa \in \{0, 0.05, 0.1, 0.2, 0.4\}$ an $8000$-member ensemble is
integrated forward to $\tau = 20$, and from the same stored turning state
two return legs are run.  One flips only the system momenta, which is
what the observer can do.  The other flips the bath momenta as well,
which is the control.  The witnesses are the hidden-state recovery defect
of each return, the retrace defect of the observed entropy curve, and the
plug-in mutual information between the binned observed coordinate and the
binned bath energy.

\subsection{Result}

\begin{figure}[t]
\begin{tikzpicture}
\begin{groupplot}[group style={group size=1 by 2, vertical sep=34pt},
  width=\columnwidth]
\nextgroupplot[paperaxis, height=0.5\columnwidth,
  ymode=log, ymin=1e-16, ymax=50,
  xmin=-0.02, xmax=0.42,
  ylabel={recovery defect},
  legend pos=south east]
\addplot[red!70!black, thick, mark=*, mark size=1.6pt]
  table[x=kappa, y=sysflip] {figdata/pe4_defects.dat};
\addlegendentry{system-only flip}
\addplot[blue!70!black, thick, mark=square*, mark size=1.4pt]
  table[x=kappa, y=fullflip] {figdata/pe4_defects.dat};
\addlegendentry{full flip with bath}
\nextgroupplot[paperaxis, height=0.5\columnwidth,
  xmin=-0.02, xmax=0.42, ymin=-0.3, ymax=6.8,
  xlabel={coupling $\kappa$},
  ylabel={[bits]},
  legend pos=south east]
\addplot[red!70!black, thick, mark=*, mark size=1.6pt]
  table[x=kappa, y=retrace] {figdata/pe4_defects.dat};
\addlegendentry{lost retraceability}
\addplot[blue!70!black, thick, mark=square*, mark size=1.4pt]
  table[x=kappa, y=mi] {figdata/pe4_defects.dat};
\addlegendentry{$I(Z;E_{\text{bath}})$}
\end{groupplot}
\end{tikzpicture}
\caption{Top panel.  Root-mean-square recovery defect of the hidden
system state after the round trip, for the reversal the observer can
perform and for the control that also flips the bath.  The observer's
reversal fails monotonically with coupling while the full flip recovers
to better than $5 \times 10^{-15}$ at every coupling.  Bottom panel.  The
entropy retrace defect grows to $5.87$ bits while the mutual information
between the observed coordinate and the bath energy never leaves the
estimator floor.  Data from the committed evidence record of the bath
experiment.}
\label{fig:pe4}
\end{figure}

Figure~\ref{fig:pe4} carries both findings.  At $\kappa = 0$ the cycle
reproduces the published result exactly, with recovery defect
$4.8 \times 10^{-15}$ and entropy retrace defect exactly zero.  As the
coupling grows, the reversal available to the observer fails
monotonically, with recovery defects $0.33$, $0.64$, $1.18$, and $1.89$
across the grid and a retrace failure that reaches $5.87$ bits.  The
control tells the other half of the story.  The full flip, including the
bath, recovers the hidden state to better than $5 \times 10^{-15}$ at
every coupling.  The dynamics never became irreversible.  What changed is
the reach of the observer's reversal protocol.  In this model class,
entropy production is an accounting of which degrees of freedom the
observer can flip \cite{Landauer1961, Bennett1982}, and the published
claim that observed entropy carries
no hidden loss survives intact in the enlarged space.

The second finding was not the expected one, and the runner's first
version asserted its opposite.  The natural expectation is that the
leaked information should be visible in the bath, for instance as mutual
information between the observed coordinate and the bath energy.
Measured, it is not.  Across the whole grid that mutual information stays
between $0.0086$ and $0.0162$ bits, within a factor of two of its
zero-coupling estimator floor, while the observer's lost retraceability
grows to almost six bits.  The mean bath energy is no witness either,
since without a counter-term the coupling shifts the equilibrium and the
bath energy actually falls at strong coupling.  The record of the leak
lives in system-bath microstate correlations that none of the coarse bath
functionals measured here can see.  Recovery therefore requires
microstate access, not a macroscopic bath reading.  The revised assertion
set demands this invisibility, and a coarse mutual information rising
above $0.05$ bits would now refute the reading.  This is the classical
face of a familiar quantum lesson, that what an environment records and
what a coarse observer of that environment can read are very different
things \cite{Zurek2003}.

\section{What the track establishes}
\label{sec:claims}

\begin{table*}[t]
\caption{Claim register for the completed entropy track.  Labels follow
the campaign's evidence scheme.  Proved means a theorem with a written
proof.  Exact means a closed-form computation verified in independent
implementations.  Demonstrated in model means a measured simulation with
declared parameters, append-only evidence records, and falsification bars
fixed before the run.  Measured in model, exploratory label, means the
same with the evidence record carrying the exploratory label because no
seal governs the run.}
\label{tab:claims}
\begin{ruledtabular}
\begin{tabular}{L{4.6cm} L{9.2cm} l}
Experiment & Claim & Label \\
\colrule
Fold theorem \cite{EntropyPaper} & a nondegenerate fold changes fiber
multiplicity by two and conserves the signed count & proved \\
Branch entropies \cite{EntropyPaper} & fold one bit, symmetric double
fold one and one half bits, degenerate monotone map zero & exact \\
Caustic scaling \cite{EntropyPaper} & binned caustic entropy converges to
its differential limit as the square root of the bin width & demonstrated
in model \\
Reversible cycle \cite{EntropyPaper} & observed entropy swings above two
bits with zero hidden loss and exact retrace & demonstrated in model \\
Mixing versus multiplicity (this paper) & sustained entropy growth
follows the dynamics, multiplicity follows the observer, independently &
measured in model, exploratory label \\
Consumer hierarchy (this paper) & exact conditional chain with positive
gaps pricing orientation and branch knowledge separately & measured in
model, exploratory label \\
Bath coupling (this paper) & reversal failure is set by the observer's
reach, the full flip recovers exactly at every coupling, and the leak is
invisible to coarse bath observables & measured in model, exploratory
label \\
\end{tabular}
\end{ruledtabular}
\end{table*}

Table~\ref{tab:claims} registers the claims of the completed track with
their labels.  Read together, the six experiments give the campaign's
answer to where observed entropy comes from in this model class.  The
growth of observed entropy is mixing, a property of the hidden dynamics.
The ambiguity of an observation is multiplicity, a property of the
observer's own singularities, worth an exactly computable number of bits
to resolve.  The irreversibility of observed entropy is neither of these.
It is an accounting statement about which degrees of freedom the
observer's reversal can reach, it appears at the moment ambiguity leaks
into unreachable ones, and it leaves no record that a coarse reading of
the environment can find.

Three cautions bound the reading.  These are toy Hamiltonians with a
handful of degrees of freedom, and no thermodynamic limit is taken
anywhere in this paper.  The bath experiment measures one family of
coarse bath observables, and a cleverer functional of the bath might see
more of the leak than the energy does.  And the mixing experiment's
chaotic member was diagnosed by its behavior rather than by a Lyapunov
measurement, so the label chaotic is operational.  None of these cautions
weakens the matched comparisons, which are internal to each experiment.

\section{Methods and reproducibility}
\label{sec:methods}

All instruments derive from the campaign's sealed tolerance and
instrument freezes, registered as PF0-FREEZE-001 and PEQO-FREEZE-002,
which fix numerical bars in advance and void themselves if edited.  The
three experiments of this paper are exploratory-grade instruments
validated against exact controls, the zero-coupling limit for the bath
cycle, the monotone observer for the multiplicity count, and the exact
finite sums for the hierarchy.  Every result-bearing run writes an
append-only JSON evidence record containing the complete parameter set,
seeds, outcome arrays, and a canonical SHA-256 hash of the record body,
and the records for the three experiments here are
\texttt{pe3-mixing.json}, \texttt{pe5-hierarchy.json}, and
\texttt{pe4-noise.json} in the campaign repository.  Every number and
every figure coordinate in this paper is read from those committed
records or recomputed deterministically from committed instrument code by
a committed script, with no free parameters.  One instrument correction
made during this work is recorded in the repository history.  The
degenerate initial ensemble of the bath cycle initially sat exactly on a
bin edge, where round-trip noise at the $10^{-16}$ level split one bin
into two and produced a spurious retrace defect, and the binning grid was
shifted by half a bin before any evidence was recorded.

\begin{thebibliography}{9}

\bibitem{EntropyPaper}
A.~H.~Bond,
Entropy across singular projections: fiber multiplicity, caustics, and
observer-relative information,
Zenodo (2026), doi:10.5281/zenodo.21789011.

\bibitem{QuantumPaper}
A.~H.~Bond,
Consumer-relative quantum distinguishability,
Zenodo (2026), doi:10.5281/zenodo.21789046.

\bibitem{CoverThomas2006}
T.~M.~Cover and J.~A.~Thomas,
\emph{Elements of Information Theory}, 2nd ed.\
(Wiley-Interscience, Hoboken, 2006).

\bibitem{CaldeiraLeggett1983}
A.~O.~Caldeira and A.~J.~Leggett,
Path integral approach to quantum Brownian motion,
Physica A \textbf{121}, 587 (1983).

\bibitem{Zurek2003}
W.~H.~Zurek,
Decoherence, einselection, and the quantum origins of the classical,
Rev.\ Mod.\ Phys.\ \textbf{75}, 715 (2003).

\bibitem{Landauer1961}
R.~Landauer,
Irreversibility and heat generation in the computing process,
IBM J.\ Res.\ Dev.\ \textbf{5}, 183 (1961).

\bibitem{Bennett1982}
C.~H.~Bennett,
The thermodynamics of computation, a review,
Int.\ J.\ Theor.\ Phys.\ \textbf{21}, 905 (1982).

\end{thebibliography}

\end{document}
